\documentclass[12pt,a4paper]{letter}

\usepackage[margin=2.5cm]{geometry}
\usepackage{booktabs}
\usepackage{hyperref}
\usepackage{parskip}
\usepackage{microtype}

\signature{Robin}
\address{}
\date{13 May 2026}

\begin{document}

\begin{letter}{Professor Neil Ghani}

\opening{Hi Neil,}

I've been reading through some recent work on multi-agent LLM systems and thought you might find it interesting, given our conversations about hierarchical composition and coordination. Three pieces stood out:

\bigskip
\noindent\textbf{1.\quad SciAgent: A Unified Multi-Agent System for Generalistic Scientific Reasoning}\\
\textit{arXiv:2511.08151} \quad\textemdash\quad \textit{Note: currently withdrawn pending revision}

\smallskip
SciAgent targets cross-disciplinary scientific reasoning at expert level. A \emph{Coordinator Agent} classifies the domain and difficulty of each problem, then dynamically assembles a \emph{Worker System} of sub-agents handling symbolic deduction, conceptual modelling, numerical computation, and verification. The claimed results were striking\textemdash gold-medalist level on the IMO, IPhO, IChO, and Humanity's Last Exam\textemdash but the paper has since been withdrawn by the authors pending further expert evaluation. I include it here as an interesting architecture sketch, not as a reliable empirical claim.

\bigskip
\noindent\textbf{2.\quad AgentOrchestra: A Hierarchical Multi-Agent Framework for General-Purpose Task Solving}\\
\textit{arXiv:2506.12508}

\smallskip
AgentOrchestra is a two-tier framework targeting general task-solving rather than a specific domain. A \emph{Planning Agent} decomposes tasks and routes work to three specialists:
\begin{itemize}
  \item \emph{Deep Researcher} --- broad web information gathering
  \item \emph{Browser Use Agent} --- precise web interaction and extraction
  \item \emph{Deep Analyzer} --- multimodal reasoning over retrieved content
\end{itemize}

The design principles (extensibility, multimodality, modularity, coordination) are deliberately architectural rather than empirical. Results on standard benchmarks:

\medskip
\begin{center}
\begin{tabular}{lcc}
\toprule
\textbf{Benchmark} & \textbf{AgentOrchestra} & \textbf{Previous SOTA} \\
\midrule
SimpleQA  & 95.3\%  & 93.9\% (Perplexity Deep Research) \\
GAIA      & 82.42\% & 77.58\% (AWorld) \\
HLE       & 25.9\%  & 20.3\% (o3) \\
\bottomrule
\end{tabular}
\end{center}

\medskip
The HLE result in particular is interesting: 25.9\% vs 20.3\% for o3 is a meaningful jump on a benchmark designed to be extremely difficult.

\bigskip
\noindent\textbf{3.\quad Building a Multi-Agent Research System}\\
\textit{Anthropic Engineering Blog, 2025}

\smallskip
Anthropic describe their own internal multi-agent research assistant. The architecture is an orchestrator--worker pattern: a lead \texttt{claude-opus-4} agent builds a dynamic research strategy and spawns 3--5 parallel \texttt{claude-sonnet-4} subagents, each with its own context window, doing web search and extended thinking. A separate citation agent handles attribution post-hoc.

Key empirical findings:
\begin{itemize}
  \item Multi-agent (Opus lead + Sonnet workers) outperforms single-agent Opus by \textbf{90.2\%} on internal research evaluations
  \item Parallel execution cut research time by up to \textbf{90\%} on complex queries
  \item Token usage explains \textbf{80\%} of performance variance in browsing tasks
  \item Improved tool descriptions alone reduced task completion time by \textbf{40\%}
\end{itemize}

The last point about token usage is the mechanistically interesting one: separate context windows per subagent function as parallel working memory. The performance gain is not purely about parallelism\textemdash it is about avoiding the degradation that comes from a single agent's context window filling up on a long multi-step task.

\bigskip
\noindent\textbf{Cross-cutting observation}

All three systems land on the same architectural form: a coordinator plus specialised workers, assembled hierarchically. What varies is what the coordinator \emph{does}:
\begin{itemize}
  \item \emph{SciAgent}: classifies domain $\to$ selects reasoning modalities
  \item \emph{AgentOrchestra}: decomposes task $\to$ routes to tool-appropriate specialists
  \item \emph{Anthropic}: builds a dynamic research strategy $\to$ spawns parallel explorers without pre-committing to structure
\end{itemize}

There is something here that feels adjacent to what we have been thinking about with categorical composition: the coordinator is essentially choosing a \emph{morphism} in some category of task-decompositions, and the specialists are the component objects. The hierarchical structure is not just engineering convenience\textemdash it is doing real semantic work in determining what kind of reasoning gets applied where.

I am curious whether you see a clean categorical account of the coordination problem. Seems like it could connect to your work on containers or polynomial functors.

\closing{Best,}

\end{letter}
\end{document}
